\section{Discussion and Limitations}\label{sec:discussion}

The current evidence supports the feasibility of combining shape-specific
references, matching simulation assets, and one shared controller at the
scale of the constructed corpus. It does not yet identify the best
architecture or establish a causal learning benefit from refinement.
Those conclusions depend on the corrected comparisons, not on the
pre-correction development curves.

The dataset is derived from existing motion capture, captions, body models,
and HUMOS generation. The 2,681,728 variants represent combinations of
20,951 clip identities and 128 configurations, not equally many independent
human examples. Mirrored clips are also not independent captures.
Caption alignment, left/right semantics, body-specific wording, and
redistribution permissions require further auditing before a benchmark
release or language-grounded performance claim.

Refinement changes kinematics rather than solving the complete equations
of motion. Ground clearance of collision proxies does not guarantee
balanced support, plausible contact forces, or realistic joint torques.
Similarly, estimated asset masses and mass-scaled actuators are modeling
choices rather than a calibrated description of human physiology. The
flat-ground task excludes much environmental interaction, and text-based
filtering cannot ensure that every remaining clip satisfies that scope.

Generalization is limited by the split and evaluation design. Original and
mirrored variants stay together, but the current split is not demonstrated
to be subject- or source-recording-disjoint. Validation reuses the trained
body set and has been monitored during development. Its rotating panel
covers only one shape per clip at a time, and long clips have a capped
evaluation horizon. Fixed-pair test evaluation and the unseen-morphology
experiment in Section~\ref{sec:unseen_morphology} remain necessary for
broader claims, and any morphology-generalization statement is confined to
body shapes inside the trained coefficient range rather than arbitrary
bodies.

The success threshold is permissive relative to fine-grained pose accuracy.
A high success rate should be read together with continuous errors,
smoothness, and qualitative rollouts. The logged normalized-jerk measure
also has implementation-specific windows and regularizers; it is not a
standalone certificate of physical plausibility. Raw policy-action
increments must not be interpreted as joint-angle or torque measurements.

Cost constraints motivated small-scale development and streamed full-scale
training, and introduced memory-related resumes and configuration changes on
the full-scale run. The 150-clip grid is single-seed: the architecture
conclusion is drawn from the jerk difference, which is a large effect that
recurs across both reference versions, and not from the success-rate
differences, which are within run-to-run noise and are reported as
non-discriminating. Seed repeats would sharpen the latter but are not
expected to change the former. These costs and choices should be reported
transparently; they do not by themselves establish better compute efficiency
than another method.

We explored neutral-body pretraining followed by multi-shape transfer and
mixed canonical/HUMOS source-weighted rewards during development, but did
not retain either as the main approach. Their historical measurements are
not used as controlled negative results, particularly where the earlier
evaluation protocol differed. The present method is direct multi-shape
training. Future work may revisit these alternatives after the current
comparisons and dataset audits are complete.
